\chapter{Against Madoka}
Names have power. 

There is an interesting pattern to the people Homura draws into her labyrinth: the magical girls make sense, as the main figures in her life and, at least in the most recent timeline, her teammates. Madoka's family are slightly more of a stretch, but they are people important to Madoka and therefore to Homura. Still more of a stretch are Hitomi and Kyousuke, but again, Hitomi is important to Madoka and Kyousuke is important to Hitomi, so it's not entirely unreasonable. But what possible reason could she have to bring in Kazuko (the homeroom teacher) and Nakazawa (an apparently random classmate)? And she does draw them both in—Nakazawa and the other magical girls are the only people seen to have normal faces when Homura begins doubting the reality of the people around her in math class, and both Nakazawa and Kazuko are shown unconscious on couches when the labyrinth is finally broken.

The answer lies in folklore: a witch who knows a person's name can use it in workings of magic against that person \cite{ref120}. Consider again who Homura brings into the labyrinth, and then consider the series as a whole; the list of people Homura brings into the barrier is also the list of people whose names are said or onscreen during the series and which Homura might reasonably know. (Admittedly, Madoka's father's name is not spoken onscreen, but Madoka calls him Papa, and as there is no other living father in the series, "Papa" functions well enough as a name.) 

Names have power because, in magical logic (which is, by and large, narrative logic), there is no signifier-signified distinction. The name is, in some sense, the thing named, and so to manipulate the name is to manipulate the thing \cite{ref121}. It follows, then, that if two things have the same name they must therefore be in some sense the same, that one can stand in for the other.

All of which is a roundabout way of saying that, when Homura grabs Madoka's arms and tears Madoka-the-girl out of Madoka-the-abstraction, it is an act of rebellion not just of Homura against Madoka, but of Rebellion against Madoka Magica.

And why shouldn't the film rebel against the series? Once, if a person wished to tell stories, they got up and told stories. Spoken aloud, these stories were ephemeral, changing with every telling. There were traditions, to be sure, but storytellers could be confident that their creative departures would not be seen as errors or betrayals but as the embellishments of a virtuoso performance.

Mass literacy struck a mortal blow against this form of storytelling, and radio, film, and television finished the job. This kind of storytelling lives on (as no art form ever really dies), but only as a curiosity, something to gawk at when attending a Renaissance festival or take your children to at the public library. Mostly, when we want a story, we reach for a packaged one, a book or a DVD.

This creates a challenge when an author wants to tell a cycle or series of stories, reusing the same characters or setting. The author wishes to explore and create, and in the age of oral tales was free to do so—no one particularly expected that the tales of Renard the Fox must be consistent with one another or complained, "Hey, when he seduced Leda, Zeus was a swan, how come he's a golden shower now?" After all, if the story of Leda can change with every telling, why expect it to still be the same when you hear a completely different story? 

Oral tales are living, growing, changing things. By contrast, a written or filmed tale is dead, nailed to the page or screen, unable to change or grow, fixed permanently as it was in a single telling. The audience is permitted to change and grow, so that their perspective on the tale can alter with time, but the actual creator of the tale is denied that growth. Even when it comes to crafting a sequel, audiences—"geek" or "cult" audiences notoriously so—demand continuity, which is to say they demand fealty to the tyrannical reign of dead stories. It is a wonder that more creators don't rebel! 

So Rebellion pays lip service to the series. All the events of the series clearly happened here and are given what the continuity-obsessed consider "respect," which is to say the letter of the law "Thou shalt not contradict the events of earlier entries" is slavishly obeyed. Even the structure of the film apes the structure of the show: it splits neatly into three parts, the first of which establishes a pretense of being a "normal" magical girl show that abruptly falls apart in a violent confrontation with Mami. The second (which, admittedly, has a stronger overlap with the first than in the series) then follows a magical girl as she slowly comes to the realization that she is what she fights against, and has been a witch from the start. Finally the third involves a tremendous battle against a city-scale witch, after which reality is rewritten and a new order established.

However, where the series followed Madoka, the film follows Homura, and therein all the difference lies. Madoka is a patient, careful, but very optimistic character—she waits until the very end of the series to act, but when she does so, it is decisively, and with every intent of ending what she sees as the primary problem of her universe absolutely and with finality. Homura is cynical, headstrong, and confrontational; she flings herself into conflict after conflict, until finally her own mirror of Madoka's actions is to create a world in which Homura's primary problem—Madoka's penchant for self-sacrifice—must be dealt with continually and continuously.

This does not necessarily imply that, for most of the film, Homura is in conscious rebellion against Madoka's order. Homura is initially positioned, just as in the series, as the one questioning and disrupting the status quo, true, but that status quo (as always, represented and defended by Mami) is Homura's own dream-realm. Homura is trying to return the state of the world to what she remembers, which is to say the world of the series. She only begins to rebel intentionally after Madoka tells her that, to her, being separated from her loved ones is tremendously painful—in other words, after Homura realizes that Madoka's self-sacrifice entailed actual sacrifice. To Homura, of course, the sacrifice of Madoka is unthinkable and unforgivable, even if it is Madoka herself performing the sacrifice.

Even then, however, Homura does not act on her desire to undo Madoka's sacrifice until very late in the movie, because up until that point she has no opportunity to do so. The character who is actually in rebellion against Madoka, and therefore against Madoka Magica, for the majority of the film is Kyubey, who has orchestrated the entire situation in an attempt to usurp control of the Law of Cycles and bring back witches. It is worth remembering here that in many respects Kyubey is an (unusually unflattering) authorial stand-in, and as such it makes sense that his rebellion against Madoka is the creators' Rebellion against Madoka Magica.

Kyubey's rebellion, however, is unsurprising—he is, after all, the villain of the series, and an unrepentant villain who is still around in the sequel can be assumed to at least try to resume their villainous role. Homura, by contrast, is spectacularly, obsessively loyal to Madoka, and so the film takes pains to meticulously lay out all the elements of her rebellion: She has motivation, in the form of her conversation among the flowers with Madoka and realization that she "never should have allowed" Madoka to sacrifice herself. She has inspiration, when Kyubey reveals that Madoka can choose to re-enter the world after all, and Sayaka reveals that Madoka's Buddha-nature, her memories and powers as the Law of Cycles, can be held in storage by another. And she has opportunity, when Madoka descends to take her life and prevent her from becoming a witch in the "real world"—as Kyubey says, that which can be perceived can be interfered with.

And so Homura rises as a devil-figure, tearing "God" from her heaven and bringing her down into the world. She is the ultimate bad girl, identified by Paradise Lost-quoting graffiti and Nietzsche-chanting, tomato-throwing familiars as Satan herself. She has claimed the labels "demon," "evil," and "enemy" for herself, and made clear that she plans to act them out—which brings us to yet another rebellion...
